\chapter{Un bilancio dell'ignoto}
\label{cap:bilancio}

I capitoli precedenti hanno costruito tre strumenti distinti per avvicinarsi alla domanda centrale di questo libro. Il primo è tafonomico: una mappa di ciò che il tempo, il clima e le azioni umane cancellano o conservano in modo selettivo e prevedibile. Il secondo è critico: un'analisi delle lacune della ricerca, delle zone bianche che non riflettono l'assenza di tracce ma l'assenza di chi le cercasse. Il terzo è metodologico: i criteri che, applicati alla morfologia del territorio, identificano i tipi di luoghi più promettenti per ciascun periodo preistorico. Questo capitolo li mette insieme in due tavole sintetiche, analizza cosa dicono quando letti in parallelo, e poi indica le lacune osservative su larga scala che nessuna delle tre analisi precedenti ha ancora esaurito.

% ============================================================
\section{Come leggere le due tavole}
% ============================================================

Le tavole che seguono usano tre simboli. Il simbolo $\bullet$ indica presenza documentata: l'arte rupestre o la costruzione di quel tipo è attestata in Veneto o in un'area immediatamente comparabile per geografia, substrato e periodo. Il simbolo $\circ$ indica presenza plausibile, non documentata in Veneto ma attesa per analogia con contesti alpini simili in cui è stata cercata con metodi adeguati. Il trattino $-$ indica assenza o improbabilità nelle condizioni del Veneto, per ragioni tafonomiche, climatiche o legate alle forme di occupazione del territorio nell'epoca considerata.

La colonna \emph{costruzioni} riunisce categorie materialmente eterogenee: megaliti e stele menhir nel Neolitico e nell'Eneolitico; dossi artificiali e terrapieni nelle pianure dell'età del Bronzo; castellieri e strutture difensive sulle creste dell'età del Bronzo e del Ferro; strutture con funzioni rituali o calendariali di cui restano tracce indirette. Le accomuna un'unica caratteristica: sono modificazioni permanenti e intenzionali del paesaggio fisico, distinte dall'arte rupestre in senso stretto ma espressione dello stesso impulso di marcare il territorio in modo durevole. In Veneto alcune di queste categorie sono documentate; altre sono quasi assenti dalla letteratura non perché manchino, ma perché non sono mai state sistematicamente cercate.

La seconda tavola include le epoche prevalenti per ciascun tipo di luogo, poiché la pertinenza di un sito non dipende solo dalla sua morfologia ma dal periodo in cui era frequentato con l'intensità sufficiente a produrre arte o strutture permanenti. Un valico a duemila metri è un candidato plausibile per il Mesolitico e per l'Eneolitico; lo è molto meno per il Paleolitico superiore, quando quella quota era ghiacciata o periglaciale.

% ============================================================
\section{Tavola 1: quali tipi di arte, in quale epoca}
% ============================================================

\begin{table}[h]
\small
\setlength{\extrarowheight}{3pt}
\caption{Presenza attesa per tipo di arte e per epoca nel territorio veneto e nelle aree immediatamente comparabili. Date in migliaia di anni prima dell'anno zero astronomico. $\bullet$~=~documentato;\enspace $\circ$~=~plausibile per analogia;\enspace $-$~=~assente o molto improbabile.}
\label{tab:bilancio_tipi}
\begin{tabularx}{\textwidth}{>{\raggedright\arraybackslash}X
                              >{\centering\arraybackslash}p{2.0cm}
                              >{\centering\arraybackslash}p{2.0cm}
                              >{\centering\arraybackslash}p{2.4cm}}
\toprule
\textbf{Epoca} & \textbf{Incisioni} & \textbf{Pitture} & \textbf{Costruzioni} \\
\midrule
\rowcolor{blue!8}
Paleolitico superiore (38.000--11.500 a.p.a.z.) & $\circ$ & $\bullet$ & $-$ \\
Mesolitico (11.500--7.500 a.p.a.z.) & $\circ$ & $\circ$ & $-$ \\
\rowcolor{blue!8}
Neolitico (7.500--5.250 a.p.a.z.) & $\circ$ & $-$ & $\circ$ \\
Eneolitico (5.250--4.250 a.p.a.z.) & $\circ$ & $-$ & $\bullet$ \\
\rowcolor{blue!8}
Età del Bronzo (4.250--2.750 a.p.a.z.) & $\bullet$ & $-$ & $\bullet$ \\
Età del Ferro e periodo venetico (2.750--0 a.p.a.z.) & $\bullet$ & $-$ & $\bullet$ \\
\bottomrule
\end{tabularx}
\end{table}

% ============================================================
\section{Tavola 2: quali luoghi, quali tipi di arte, quali epoche}
% ============================================================

\small
\setlength{\extrarowheight}{3pt}
\begin{longtable}{>{\raggedright\arraybackslash}p{3.0cm}
                   >{\centering\arraybackslash}p{1.5cm}
                   >{\centering\arraybackslash}p{1.5cm}
                   >{\centering\arraybackslash}p{1.7cm}
                   >{\raggedright\arraybackslash}p{2.8cm}}

\caption{Presenza attesa per tipo di luogo e tipo di arte, con le epoche prevalenti. I giudizi tengono conto della tafonomia, del clima e delle forme di occupazione del territorio per ciascun periodo. Stessi simboli della tavola precedente.\label{tab:bilancio_luoghi}} \\

\toprule
\textbf{Tipo di luogo} & \textbf{Incis.} & \textbf{Pitt.} & \textbf{Costr.} & \textbf{Epoche prevalenti} \\
\midrule
\endfirsthead

\multicolumn{5}{l}{\small\textit{Tavola~\ref{tab:bilancio_luoghi} -- continua dalla pagina precedente}} \\[4pt]
\toprule
\textbf{Tipo di luogo} & \textbf{Incis.} & \textbf{Pitt.} & \textbf{Costr.} & \textbf{Epoche prevalenti} \\
\midrule
\endhead

\midrule
\multicolumn{5}{r}{\small\textit{continua nella pagina seguente}} \\
\endfoot

\bottomrule
\endlastfoot

\rowcolor{blue!8}
Grotte (interni) &
  $\circ$ & $\bullet$ & $-$ &
  Paleolitico, Mesolitico \\[2pt]

Ripari rocciosi semi-aperti &
  $\circ$ & $\bullet$ & $-$ &
  Paleolitico, Mesolitico \\[2pt]

\rowcolor{blue!8}
Punti di avvistamento (speroni e rilievi con visibilità multipla) &
  $\circ$ & $-$ & $\circ$ &
  Mesolitico, Bronzo \\[2pt]

Selle e valichi &
  $\circ$ & $-$ & $\circ$ &
  Mesolitico, Eneolitico, Bronzo \\[2pt]

\rowcolor{blue!8}
Pareti verticali a quota intermedia (400--1.200 m, con overhang) &
  $\bullet$ & $-$ & $-$ &
  Bronzo, Ferro \\[2pt]

Radure d'alta quota non abitate &
  $\circ$ & $-$ & $\circ$ &
  Neolitico, Bronzo \\[2pt]

\rowcolor{blue!8}
Falde dei gruppi montagnosi (zona di transizione bosco-quota) &
  $\circ$ & $-$ & $\circ$ &
  Mesolitico, Neolitico, Bronzo \\[2pt]

Crode dolomitiche d'alta quota (oltre 1.800 m) &
  $\circ$ & $-$ & $\circ$ &
  Bronzo, Ferro (uso cultuale) \\[2pt]

\rowcolor{blue!8}
Fondovalle montani (Val Belluna, Zoldano, Agordino, Piave alta) &
  $\circ$ & $-$ & $\circ$ &
  Neolitico, Bronzo, Ferro \\[2pt]

Colli isolati di pianura (Euganei, Berici) &
  $\circ$ & $-$ & $\circ$ &
  Paleolitico, Mesolitico, Neolitico \\[2pt]

\rowcolor{blue!8}
Dossi e terrazze fluviali di pianura &
  $-$ & $-$ & $\bullet$ &
  Neolitico, Bronzo, Ferro \\[2pt]

\end{longtable}
\normalsize

% ============================================================
\section{Cosa dicono le tavole}
% ============================================================

Lette insieme, le due tavole producono alcune osservazioni che non emergono separatamente dall'analisi tafonomica o da quella metodologica.

La prima: le pitture sono il tipo di arte rupestre più fragile all'aperto e al tempo stesso il meglio documentato in Veneto per il Paleolitico superiore. Fumane e Dalmeri sono pitture, non incisioni. Non è casuale: le pitture su parete interna di grotta, in penombra e protette dall'umidità costante, sopravvivono meglio di qualsiasi altra categoria in ambiente calcareo. La loro assenza dal mondo esterno è quasi totale per ragioni tafonomiche prevedibili; la loro presenza nelle grotte della Lessinia e del Veneto orientale è sottodocumentata perché non è mai stata sistematicamente cercata con metodi ottici avanzati.

La seconda: le incisioni all'aperto sono documentate solo per il Bronzo e il Ferro, ma plausibili per quasi tutte le epoche precedenti. Il divario non è nella produzione originaria, ma nella sopravvivenza. Un'incisione su calcare esposto alle precipitazioni perde leggibilità in poche centinaia di anni; un'incisione su calcare protetto da un overhang può durare millenni. La differenza tra un $\bullet$ e un $\circ$ nella colonna delle incisioni riflette quasi sempre la presenza o l'assenza di protezione morfologica, non la presenza o l'assenza dell'attività umana.

La terza: le costruzioni sono il tipo di traccia più resistente al tempo quando sono in pietra o in terra compattata, ma la più invisibile quando erano in legno o in materiale organico. I castellieri dell'età del Bronzo sopravvivono perché erano strutture in pietra su creste già rocciose. I probabili punti di sosta mesolitici non lasciano nulla di visibile in superficie. Il divario tra i $\bullet$ delle costruzioni documentate e i $\circ$ di quelle plausibili nelle epoche più antiche non misura l'intensità dell'occupazione umana ma la sua modalità: stabile e sedentaria contro mobile e stagionale.

La quarta, e forse la più importante per il ragionamento di questo libro: le righe con più $\circ$ nella seconda tavola non sono le meno promettenti. Sono le meno cercate. Un $\circ$ significa che l'analogia con aree comparabili dice che ci dovrebbe essere qualcosa, e che la ragione per cui non è documentato è che nessuno si è fermato a guardare con metodi adeguati.

% ============================================================
\section{Le lacune osservative su larga scala}
% ============================================================

Il capitolo sul research bias ha già analizzato le lacune legate ai siti noti: i meccanismi che producono mappe della conoscenza disomogenee, dove la presenza di un museo, di un naturalista ottocentesco o di una cava produce concentrazioni di segnalazioni che nulla hanno a che fare con la distribuzione reale dell'arte preistorica. Quello che rimane fuori da quella analisi sono le lacune geografiche strutturali: intere aree del territorio veneto e delle zone contigue che non figurano in alcun programma di ricerca sistematica per arte rupestre preistorica. Non si tratta di singole pareti o di piccole creste, ma di masse montuose, fasce altitudinali e bacini vallivi di scala tale da costituire, nella loro assenza dalla letteratura, un vuoto obiettivo.

\subsection*{Le crode del Cadore e del Comelico}

I principali gruppi dolomitici del Cadore e del Comelico, la Marmolada, le Pale di San Martino, le Pale di San Lucano, la Civetta, l'Antelao, il Pelmo, l'Averau, il Nuvolau, le Tofane, il Sorapiss, le Marmarole, il Monte Schiara, le Vette Feltrine e tutti i gruppi minori intermedi, non risultano oggetto di alcuna prospezione sistematica per arte rupestre preistorica. I pochi segnali disponibili confermano che la ricerca c'è stata appena abbozzata: Fabio Cavulli ha presentato nel 2025 a Pieve di Cadore i risultati preliminari di cerchi incisi al compasso su superfici dolomitiche, un indizio che queste pareti potrebbero riservare sorprese.

Ma c'è un secondo motivo per non escludere le alte quote, un motivo che va oltre la caccia e i campi base: il culto delle montagne. Constanza Ceruti, l'archeologa argentina che ha ritrovato mummie sacrificali sulle cime andine a oltre sei mila metri di quota, ha documentato in un recente lavoro sull'arco alpino come la sacralità delle cime alpine, comprese le Dolomiti italo-austriache, sia rimasta largamente ignorata dalla discussione accademica contemporanea \citep{ceruti2025}. Il riferimento andino non è decorativo: Ceruti e Reinhard hanno dimostrato che il culto delle montagne è documentato in modo indipendente su quattro continenti, dalle Ande all'Himalaya, dalle Alpi alle catene messicane \citep{reinhard2010}, il che suggerisce che risponda a qualcosa di radicato nell'esperienza umana del paesaggio d'alta quota. Non sarebbe la prima volta che le Alpi conservano tracce di questa relazione: nel 2024, quando il ghiacciaio del Tresero in Alta Valtellina si è ritirato abbastanza da esporre la roccia sottostante, sono comparse le incisioni rupestri più alte d'Europa, a tremila metri di quota, probabilmente parte di un santuario rupestre dell'età del Bronzo che il ghiaccio aveva coperto per millenni \citep{tresero2024}. Il ritiro glaciale accelerato degli ultimi decenni sta rendendo visibili superfici che erano inaccessibili durante tutta la storia della ricerca moderna. Le Dolomiti venete hanno avuto ghiacciai propri.

\subsection*{Le falde dei gruppi montagnosi}

I versanti inferiori e le fasce di transizione tra bosco e quota dei principali massicci veneti, Baldo, Carega, Pasubio, Recoaro, Tonezza, le zone dell'Altopiano di Asiago non interessate dai combattimenti e il Grappa, sono aree frequentate dalla vita umana moderna ma non sistematicamente cercate per arte preistorica. Esistono già ritrovamenti in alcune di queste zone: il Baldo, Asiago e Folgaria hanno restituito materiali che attestano frequentazione preistorica. Questi ritrovamenti non sono la conclusione dell'indagine ma il punto di partenza: dimostrano presenza umana in queste fasce altitudinali, il che rende plausibile la presenza di arte rupestre nei tipi di luogo descritti nel capitolo precedente. Le falde a quote tra i quattrocento e i milleduecento metri sono precisamente la fascia in cui il capitolo sulla metodologia individua le pareti verticali con overhang dell'età del Bronzo e i ripari semi-aperti del Mesolitico.

\subsection*{I fondovalle montani non indagati}

La Val Belluna, la valle del Piave nella sua sezione alta, lo Zoldano e l'Agordino sono pressoché assenti dalla letteratura sull'arte rupestre preistorica veneta. Questi fondovalle non sono ambienti sfavorevoli: sono corridoi naturali di transito tra la pianura e le quote alpine, con substrati calcarei nelle pareti laterali delle valli e microclimi stabili che avrebbero favorito sia la frequentazione sia la conservazione delle tracce. La loro assenza dalla letteratura è una lacuna di ricerca, non un risultato.

\subsection*{Le radure non abitate}

Le radure d'alta quota sono ambienti trasformati dall'azione umana, ma non necessariamente da un'azione recente. La maggior parte dei pianori erbosi delle Prealpi venete è stata disboscata a partire dal Medioevo, e le strutture pastorali visibili in superficie, malghe, tabià, muretti a secco, datano agli ultimi cinque-otto secoli. Sotto queste strutture e nei loro margini possono esistere tracce di frequentazione molto più antica: coppelle su pianori rocciosi calcarei, superfici sub-pianeggianti con incisioni, depositi di materiale in anfratti protetti. Le radure attualmente in uso per la pastorizia sono difficili da esplorare sistematicamente; quelle abbandonate, dove la vegetazione ha ripreso il sopravvento dopo la cessazione dell'uso, sono quasi totalmente ignorate dalla ricerca. Sono anche le meno disturbate dalle attività recenti.

\subsection*{I Colli Euganei}

I Colli Euganei hanno una geologia anomala rispetto al resto del Veneto: sono un massiccio di rocce magmatiche intrusive ed effusive, trachiti, fonoliti e rioliti di età eocenica, che emergono isolati dalla pianura fino a 601 metri. Come già discusso nel capitolo sulla storia del territorio, questi colli hanno funzionato per tutta la preistoria come rifugi ecologici con microclima mite, acqua affidabile e substrato duro favorevole alla scheggiatura. Ciò che non è mai stato sistematicamente cercato è l'arte rupestre su roccia vulcanica. Gli affioramenti trachitici hanno superfici dure e resistenti, molto più favorevoli all'incisione duratura rispetto al calcare. La presenza di minerali cupriferi e ferrosi nei Colli Euganei rende geologicamente plausibile che comunità eneolitiche e dell'età del Bronzo abbiano riconosciuto e frequentato questi colli non solo come rifugio ma come fonte di materie prime: e dove c'era estrazione e frequentazione intensiva, c'era anche la condizione che altrove ha prodotto arte rupestre. Non esistono ancora prove, ma esistono i presupposti perché valga la pena cercare.

% ============================================================
\section{Quello che potremmo trovare}
% ============================================================

Mettendo insieme la tafonomia, le lacune della ricerca e la metodologia di prospezione, si può dire qualcosa di ragionevole sulla natura di ciò che potrebbe emergere da una ricerca sistematica del territorio veneto, senza cedere né all'ottimismo facile né al pessimismo ingiustificato.

I ripari e le grotte sono il tipo di contesto dove la probabilità di trovare qualcosa è concretamente più alta. Le pitture del Paleolitico superiore nelle grotte della Lessinia e dei gruppi prealpini non sono mai state cercate con metodi ottici avanzati. I ripari rocciosi delle falde prealpine, tra i quattrocento e gli ottocento metri, potrebbero conservare incisioni su pareti protette dall'overhang che la vegetazione e la patina superficiale rendono invisibili alla ricognizione ordinaria. Questi contesti hanno la caratteristica tafonomica più favorevole: sono chiusi o semi-chiusi, protetti dalle precipitazioni, e hanno subito variazioni di temperatura e umidità meno estreme di qualsiasi superficie all'aperto.

L'arte rupestre su pareti aperte è meno probabile, per le ragioni discusse nel capitolo sulla tafonomia: il calcare delle Prealpi venete, esposto alle precipitazioni e ai cicli di gelo e disgelo per diecimila anni, ha perso la maggior parte delle superfici che in origine potrebbero aver portato incisioni superficiali. L'eccezione è rappresentata dalle pareti con overhang pronunciato, dove la protezione è stata sufficiente a rallentare l'erosione: e queste pareti sono precisamente il target metodologico del capitolo precedente.

Le crode d'alta quota restano l'incognita più aperta, e forse la più affascinante. Non per la probabilità statistica di un ritrovamento, che è bassa, ma perché un ritrovamento in quel contesto avrebbe un significato diverso da tutti gli altri: non un sito di caccia, non un campo base stagionale, ma qualcosa che richiederebbe una spiegazione diversa. Il ghiacciaio del Tresero ha mostrato che questa categoria esiste nelle Alpi italiane. Il ritiro glaciale in corso sta aprendo superfici che nessun ricercatore ha mai potuto vedere. Le Dolomiti venete hanno ghiacciai che si ritirano. Non sappiamo cosa nascondono.

Basta trovarne uno. Un riparo con pitture mai documentate, una grotta con incisioni che nessuno aveva cercato, un segno su una croda a quota elevata che non ha spiegazioni immediate. Cambierebbe tutto ciò che ora riteniamo di sapere sul silenzio di queste montagne.
